%! Justifying a Claim Based on a Confidence Interval for a Population Proportion — Unit 6, Lesson 6.3 — Group Activity
\documentclass[10pt]{article}
\usepackage{apstats-article}
\usepackage{apstats-boxes}
\newcommand{\TallMath}[1]{$\displaystyle #1\rule[-1.4em]{0pt}{3.2em}$}

\begin{document}

\pageheader{Unit 6, Lesson 6.3}{Group Activity}
\namepartnerperiod

\vspace{0.15in}

\textit{Each scenario gives you a computed confidence interval. Interpret it, evaluate the
stated claim, and answer the follow-up question.}

\begin{scenariobox}[Scenario 1 --- Flu Vaccine Uptake]{sky}
A random sample of 500 adults found $\hat p = 0.46$ received a flu vaccine last season.
A 95\% CI yields $(0.417, 0.503)$.

\begin{enumerate}[label=\alph*.]
  \item Write a complete interpretation of this interval.
        \writelines{2}

  \item A public health official claims that fewer than half of adults get the flu vaccine
        ($p < 0.50$). Does the CI support, refute, or neither support nor refute this claim?
        Explain.
        \writelines{2}

  \item How would the width of the interval change if $n = 2000$ were used instead? Would
        the claim evaluation necessarily change?
        \writelines{2}
\end{enumerate}
\end{scenariobox}

\begin{scenariobox}[Scenario 2 --- Recycling Habits]{sky}
A survey of 280 randomly selected households found 196 recycle regularly. A 90\% CI
yields $(0.652, 0.748)$.

\begin{enumerate}[label=\alph*.]
  \item Write a complete interpretation.
        \writelines{2}

  \item A city council member claims that exactly 70\% of households recycle regularly.
        Is this claim plausible based on the CI? Explain.
        \writelines{2}

  \item A second council member says ``At least two-thirds of households recycle.''
        Does the CI support this? Explain.
        \writelines{2}
\end{enumerate}
\end{scenariobox}

\begin{scenariobox}[Scenario 3 --- Social Media Use]{sky}
A random sample of 150 employees found $\hat p = 0.54$ use social media during work hours.
A 99\% CI yields $(0.437, 0.643)$.

\begin{enumerate}[label=\alph*.]
  \item A manager claims the proportion is 0.50. Is this plausible? \blank{2cm} Explain: \writelines{1}

  \item A different manager claims the proportion is 0.65. Is this plausible? \blank{2cm} Explain: \writelines{1}

  \item Why is this 99\% CI unusually wide compared to most CIs in this unit?
        \writelines{2}
\end{enumerate}
\end{scenariobox}

\begin{reflectionbox}
Based on all three scenarios, write a general rule for when a CI provides clear evidence
against a claim vs.\ when the evidence is inconclusive. Use the word ``plausible.''
\writelines{3}
\end{reflectionbox}

\end{document}
